% !TEX program = lualatex
\documentclass[
11pt, english, 
%draft, % no pictures, links, overfull hboxes indicated
toctotoc, 		      % add table of contents to the table of contents
%liststotoc, 	       % add the list of figs/tables/etc to the table of contents
%nolistspacing, 	% If the document is onehalfspacing or doublespacing, uncomment this to set spacing in lists to single
%parskip, 			% add space between paragraphs
%headsepline, 		% get a line under the header
%chapterinoneline, 	% to place the chapter title next to the number on one line
%consistentlayout, 	% to change the layout of the declaration, abstract and acknowledgements pages to match the default layout
]{style_thesis}

%-------------------------------------------------------------------------------
%   Define the LIBRARY and other packages
%-------------------------------------------------------------------------------
% actually import stzle with style=customstyle???? -> not needed
\usepackage[
    backend=biber,
    sorting=none,
    style=phys,
    articletitle=true,
    biblabel=brackets,  %e.g.[1]
    citestyle=numeric-comp, %e.g., [1, 2, 3] -> [1-3]
    maxbibnames=2,
    minbibnames=1,
]{biblatex}
\renewcommand*{\bibfont}{\small} % Optional: Make the bibliography font size smaller
\addbibresource{bib/my_bibliography.bib}

% Suppress visible DOI/URL/eprint text but keep hyperlinks on journal titles
\DeclareFieldFormat{doi}{}  % Don't print "doi: xxxx"
\DeclareFieldFormat{url}{}  % Don't print URLs
\DeclareFieldFormat{eprint}{}  % Don't print arXiv IDs
\DeclareFieldFormat{isbn}{}  % Don't print ISBN
\DeclareFieldFormat{issn}{}  % Don't print ISSN

% Added content-related packages and definitions
\usepackage{graphicx} % Required to include images
\graphicspath{/figures}
\usepackage{amsmath, amssymb, amsfonts} % Math packages
\usepackage{dsfont} % gives \mathds{1}
\usepackage{physics} % For quantum mechanics notation
% Load cleveref after hyperref and configure reference naming
\AtEndPreamble{%
    \RequirePackage[capitalise,nameinlink]{cleveref}
    % Unify naming and spell out references at sentence start
    \crefname{equation}{Eq.}{Eqs.}
    \Crefname{equation}{Equation}{Equations}
    \crefname{figure}{Fig.}{Figs.}
    \Crefname{figure}{Figure}{Figures}
    \crefname{table}{Tab.}{Tabs.}
    \Crefname{table}{Table}{Tables}
    \crefname{section}{Sec.}{Secs.}
    \Crefname{section}{Section}{Sections}
    \crefname{subsection}{Sec.}{Secs.}
    \Crefname{subsection}{Section}{Sections}
    \crefname{subsubsection}{Sec.}{Secs.}
    \Crefname{subsubsection}{Section}{Sections}
    \crefname{chapter}{Chapter}{Chapters}
    \Crefname{chapter}{Chapter}{Chapters}
    \crefname{appendix}{Appendix}{Appendices}
    \Crefname{appendix}{Appendix}{Appendices}
    % Sentence-start aware \autoref using cleveref
    % More robust version: uses spacefactor with higher threshold and vmode check
    \makeatletter
    \renewcommand{\autoref}[1]{%
        \ifvmode                           
            \Cref{#1}%
        \else
            \ifnum\spacefactor>1500\relax
                \Cref{#1}%
            \else
                \cref{#1}%
            \fi
        \fi
    }%
    % Explicit alternatives when you need them:
    % \cref{label}         - Always short form (Fig., Eq., Sec.)
    % \Cref{label}         - Always full form (Figure, Equation, Section)
    % \autoref{label}      - Automatic (this command)
    \makeatother
}
% -----------------------------------------------------------------------------
% Reference usage guide (keep consistent throughout the thesis)
% -----------------------------------------------------------------------------
% Use \autoref{<label>} in running text. It automatically chooses the correct
% name and capitalization at the start of a sentence.
%   - Example (mid-sentence): "... as shown in~\autoref{fig:my_figure}."
%   - Example (sentence start): "\autoref{eq:my_equation} shows ..."
%
% Use \cref{<label>} when you explicitly want the short form (Eq./Sec./Fig.)
% or when referencing multiple labels in one place:
%   - Multiple: "... see \cref{eq:a,eq:b,eq:c}."
%   - Range: "... see \crefrange{eq:a}{eq:c}."
%
% Avoid \ref/\eqref for numbered elements. The scheme above keeps naming and
% sentence-start capitalization consistent across the whole thesis.


%-------------------------------------------------------------------------------
%   Float placement parameters (reduce white space around figures)
%-------------------------------------------------------------------------------
% Allow more floats per page
\setcounter{totalnumber}{5}       % max floats per page (default 3)
\setcounter{topnumber}{3}         % max floats at top (default 2)
\setcounter{bottomnumber}{3}      % max floats at bottom (default 1)

% Relax fraction constraints
\renewcommand{\topfraction}{0.85}       % max fraction of page for top floats (default 0.7)
\renewcommand{\bottomfraction}{0.70}    % max fraction for bottom floats (default 0.3)
\renewcommand{\textfraction}{0.10}      % min fraction for text (default 0.2)
\renewcommand{\floatpagefraction}{0.70} % min fraction for float-only page (default 0.5)

% Reduce spacing around floats
\setlength{\floatsep}{12pt plus 2pt minus 2pt}          % space between floats (default 12pt plus 2pt minus 2pt)
\setlength{\textfloatsep}{16pt plus 2pt minus 4pt}     % space between text and floats (default 20pt plus 2pt minus 4pt)
\setlength{\intextsep}{12pt plus 2pt minus 2pt}        % space around in-text floats (default 12pt plus 2pt minus 2pt)


%\usepackage{mathpazo} % TODO Use the Palatino font by default, % ONLY WORKS WITH OLD PDFLATEX VERSION? 
\newcommand{\todoimp}[1]{\todo[inline,color=red!70]{\textbf{IMPORTANT:} #1}}
\newcommand{\todoidea}[1]{\todo[inline,color=green!60]{\textbf{IDEA:} #1}}
\newcommand{\todoeq}[1]{\todo[inline,color=blue!50]{\textbf{EQUATION:} #1}}
\newcommand{\todoref}[1]{\todo[inline,color=purple!50]{\textbf{REFERENCE:} #1}}
\newcommand{\todofix}[1]{\todo[inline,color=orange!70]{\textbf{FIX:} #1}}

%-------------------------------------------------------------------------------
%	THESIS INFORMATION
%-------------------------------------------------------------------------------

\thesistitle{2D Electron Spectroscopy in Hot Environments: A Redfield Master Equation Approach}
\supervisor{\href{https://ic1.ugr.es/members/dmanzano/home/}{Prof. Dr. Daniel \textsc{Manzano Diosdado}}}
\examiner{\href{https://sites.google.com/site/beatrizolmosphysics/}{Prof. Dr. Beatriz \textsc{Olmos Sanchez}}}
\degree{Master of Science}
\author{Leopold \textsc{Bodamer}}
\addresses{Wolfsbergallee 17b, 75177 Pforzheim}
\subject{Theoretical Atomic Physics and Synthetic Quantum Systems}
\universityone{\href{https://uni-tuebingen.de\\}{Eberhard Karls Universität Tübingen}}

\department{\href{https://uni-tuebingen.de/fakultaeten/mathematisch-naturwissenschaftliche-fakultaet/fachbereiche/physik/institute/institut-fuer-theoretische-physik/arbeitsgruppen/}{Institut für Theoretische Physik}}
\group{\href{https://uni-tuebingen.de/fakultaeten/mathematisch-naturwissenschaftliche-fakultaet/fachbereiche/physik/institute/institut-fuer-theoretische-physik/arbeitsgruppen/ag-lesanovsky/}{Theoretical Atomic Physics and Synthetic Quantum Systems}}
\faculty{\href{https://uni-tuebingen.de/fakultaeten/mathematisch-naturwissenschaftliche-fakultaet/fakultaet/}{Mathematisch-Naturwissenschaftliche Fakultät}}

\departmenttwo{\href{https://www.ugr.es/en/about/organization/entities/department-electromagnetism-and-matter-physics}{Department of Electromagnetism and Matter Physics}}
\grouptwo{\href{https://ic1.ugr.es/members/qtc/}{Quantum Thermodynamics and Computation Group}}
\facultytwo{\href{https://fciencias.ugr.es/}{Faculty of Sciences}}
\universitytwo{\href{https://www.ugr.es\\}{University of Granada}}
% NOTE could also add a faculty / adress of the universities

%\includeonly{chapters/c10_introduction} % Place this BEFORE \begin{document} to only compile specific chapter

\begin{document}
%\ The width of this document is: \the\textwidth

\frontmatter
\pagestyle{plain}

% CHAPTER 20 ONLY - STANDALONE VERSION
% Comment out title page and abstract for chapter-only compilation
% !TEX root = ../main.tex
%----------------------------------------------------------------------------------------
%	TITLE PAGE
%----------------------------------------------------------------------------------------

\begin{titlepage}
		\thispagestyle{empty}
	\begin{tikzpicture}[remember picture,overlay]
		\draw[line width=0.8pt] ([xshift=1.0cm,yshift=-1.0cm]current page.north west)
		rectangle ([xshift=-1.0cm,yshift=1.0cm]current page.south east);
	\end{tikzpicture}

	\begin{center}
		\vspace*{0.6cm}

		%\vspace*{.06\textheight}
		{\scshape\LARGE \univone\\
			\&\par
			\univtwo\\
			\par}\vspace{1.0cm} % University name

		\textsc{\Large Master Thesis}\\[0.5cm] % Thesis type

		\HRule \\[0.4cm] % Horizontal line
		{\parbox{0.9\linewidth}{\centering\huge \bfseries \ttitle\par}}\vspace{0.4cm} % Thesis title

		\HRule \\[0.9cm] % Horizontal line

		\large
		\emph{Author:}\\
		{\authorname} % Author name
		\vspace{0.7cm}

		\large
		\emph{Supervisor:} \\
		\supname \\
		\textit{from the }\\
		\grouptwoname, \depttwo\\[0.1cm]
		\factwoname\\[0.1cm]
		\emph{\univtwo}\\
		\vspace{0.7cm}

		\large
		\emph{Examiner:} \\
		\examname \\
		\textit{from the }\\
		\groupname, \deptname\\[0.1cm]
		\facname\\[0.1cm]
		\emph{\univone}\\
		\vspace{0.7cm}

		\HRule \\[0.4cm]

		\large \textit{A thesis submitted in fulfilment of the requirements\\ for the degree of \degreename}\\[0.3cm]

		\vfill
		{\large \today}\\[0.2cm] % Date
	\end{center}
\end{titlepage}

%\input{chapters/c01_declaration}
% !TEX root = ../main.tex
%-------------------------------------------------------------------------------
%	ACKNOWLEDGEMENTS
%-------------------------------------------------------------------------------

\begin{acknowledgements}
	\addchaptertocentry{\acknowledgementname}
	I would like to express my deepest gratitude to those who supported me
	throughout this journey.

	First and foremost, I sincerely thank my supervisor, \supname, for the
	continuous guidance, support, and encouragement
	throughout this work.
	The opportunity to carry out this research under his supervision and to
	gain such a valuable and unforgettable research experience made this year
	truly special for me.

	I would also like to express my sincere gratitude to my examiner,
	\examname, for making this master's thesis possible in the first place.
	By establishing contact with my supervisor, placing her trust in me, and
	supporting the organisational aspects throughout the process, she made both
	this work and my year abroad possible.
	I am deeply grateful for her support.

	Special thanks go to \'Alvaro Tejero Moyano, Rhea Alexander, Gustavo
	Mereles Menesse, Giulio Camillo da Silva, and the other members of the
	research group for their kindness, helpfulness, and the pleasant working
	atmosphere.
	Their support made this experience both productive and enjoyable.

	Finally, I would like to express my heartfelt gratitude to my girlfriend,
	Anna Schupeta, for her constant love, patience, and unwavering support throughout
	this entire process.
	She supported me tremendously through all the highs and lows, both
	emotionally and academically, and I am deeply grateful for that.
	I am also very thankful to my mother, whose outside perspective and
	thoughtful advice often provided valuable support.

\end{acknowledgements}

%-------------------------------------------------------------------------------
%	LIST OF CONTENTS/FIGURES/TABLES PAGES
%-------------------------------------------------------------------------------

\tableofcontents % Prints the main table of contents

%\listoffigures % Prints the list of figures
%\listoftables % Prints the list of tables

\clearpage
%-------------------------------------------------------------------------------
%	NOTATION AND CONVENTIONS
%-------------------------------------------------------------------------------
\subsection*{Notation and conventions}

Throughout this thesis, natural units are used with $\hbar = 1$ and $c = 1$.
Frequencies, energies, and inverse times are expressed in the same units.

The following notation conventions apply:

\begin{itemize}

\item \textbf{Vectors.}
Vectors in real space are denoted by bold symbols, e.g.\ $\boldsymbol{E}$ (electric field), $\boldsymbol{P}$ (polarisation), $\boldsymbol{k}$ (wave vector), $\boldsymbol{d}$ (transition dipole moment), and $\boldsymbol{e}$ (polarisation unit vector). 
The magnitude of a vector $\boldsymbol{d}$ is written as $d = |\boldsymbol{d}|$ when required.

\item \textbf{Operators and states.}
Operators acting on Hilbert spaces carry a hat, e.g.\ $\hat{H}$ for the Hamiltonian and $\hat{\boldsymbol{\mu}}$ for the dipole operator. 
Density operators are written as $\hat{\rho}$. 
Matrix elements in a chosen basis are denoted by plain symbols, e.g.\ $\rho_{ij} = \langle i | \hat{\rho} | j \rangle$.

\item \textbf{Quantum mechanical pictures.}
Different dynamical representations (Schr\"odinger, interaction, Heisenberg) are indicated by subscripts. 
For example, $\hat{O}_I(t)$ denotes an operator in the interaction picture. 
To avoid ambiguity with this convention, the interaction Hamiltonian is written as $\hat{H}_{\mathrm{int}}$.

\item \textbf{Superoperators.}
Linear maps acting on operators (e.g.\ propagators or Liouvillians) are written in calligraphic font, such as $\mathcal{U}$, $\mathcal{L}$. 
Their action on density operators is written explicitly as $\mathcal{L}[\hat{\rho}]$.

\item \textbf{Field projections.}
For a fixed laboratory polarisation direction $\boldsymbol{e}$, vector fields are often projected onto this axis. 
The scalar field envelope is then defined as
\[
E(t) = \boldsymbol{e} \cdot \boldsymbol{E}(t),
\]
and correspondingly $\hat{\mu} = \hat{\boldsymbol{\mu}} \cdot \boldsymbol{e}$. 
Whenever such projections are used, the underlying vector character is understood but suppressed in the notation.

\item \textbf{Complex quantities.}
The symbols $\Re$ and $\Im$ denote the real and imaginary parts of a complex quantity, respectively.

\end{itemize}

%-------------------------------------------------------------------------------
%	ABBREVIATIONS
%-------------------------------------------------------------------------------

\iffalse
	\begin{abbreviations}{ll} % Include a list of abbreviations (a table of two columns)

		\textbf{2DES} & \textbf{2D}imensional \textbf{E}lectronic \textbf{S}pectroscopy\\
		\textbf{NIR} & \textbf{N}ear \textbf{I}nfrared\\
		\textbf{IR} & \textbf{I}nfrared\\
		\textbf{UV} & \textbf{U}ltraviolet\\
		\textbf{NP} & \textbf{N}on\text{P}erturbative\\
		\textbf{FWM} & \textbf{F}our \textbf{W}ave \textbf{M}ixing\\
		\textbf{DM} & \textbf{D}ensity \textbf{M}atrix\\

	\end{abbreviations}
\fi

%-------------------------------------------------------------------------------
%	PHYSICAL CONSTANTS/OTHER DEFINITIONS
%-------------------------------------------------------------------------------

\iffalse
	\begin{constants}{lr@{${}={}$}l} % The list of physical constants is a three column table

		Constant Name & $Symbol$ & $Constant Value$ with units\\
		Speed of Light & $c$ & $1$\\
		Planck's Constant & $\hbar$ & $1$\\

	\end{constants}
\fi

%-------------------------------------------------------------------------------
%	SYMBOLS
%-------------------------------------------------------------------------------

\iffalse
	\begin{symbols}{lll} % Include a list of Symbols (a three column table)

		$a$ & distance & natural-unit length scale \\
		$P$ & polarisation amplitude & natural units \\
		%Symbol & Name & Unit \\

		\addlinespace % Gap to separate the Roman symbols from the Greek

		$\omega$ & angular frequency & inverse time (natural units) \\

	\end{symbols}
\fi


%-------------------------------------------------------------------------------
%	DEDICATION
%-------------------------------------------------------------------------------
\iffalse
	\dedicatory{For/Dedicated to/To my\ldots}
\fi


%-------------------------------------------------------------------------------
%	ABSTRACT PAGE (commented out for chapter-only version)
%-------------------------------------------------------------------------------
%\begin{abstract}
%	\addchaptertocentry{\abstractname}
%	This thesis investigates the application of two-dimensional photon-echo spectroscopy to model quantum systems, with a focus on open quantum system dynamics. A comprehensive theoretical framework is developed that combines the response-function formalism of nonlinear spectroscopy in a non-perturbative setting with master-equation approaches for open quantum systems (QuTiP). Starting from minimal reference models, including a single qubit and a four-level system formed by two coupled qubits, the methods are validated and the approach is applied to a system of $N$ qubits arranged on a cylindrical geometry inspired by the structural motifs of microtubules. The results demonstrate the capability of two-dimensional spectroscopy to probe coherence and energy-transfer dynamics in these systems, providing insights into their quantum behavior and interactions with the environment. This work lays the groundwork for future experimental and theoretical studies in the field of quantum biology and quantum information science.
%\end{abstract}

%-------------------------------------------------------------------------------
%	CONTENTS
%-------------------------------------------------------------------------------

\mainmatter
\pagestyle{thesis}

% ONLY INCLUDE CHAPTER 20: OPEN QUANTUM SYSTEMS
\include{chapters/c20_open_quantum_systems}

%-------------------------------------------------------------------------------
%	BIBLIOGRAPHY
%-------------------------------------------------------------------------------

\printbibliography

\end{document}
